% Caderno 3 — banda de confiança da resposta prevista, em ampulheta.
% Figura conceitual, sem dado. O estrangulamento em x-barra está de propósito
% mais pronunciado do que numa banda real: é ele que a figura precisa ensinar.
\documentclass[border=3pt,11pt]{standalone}
% Preâmbulo comum das figuras TikZ das apostilas.
% Carregado por todos os `standalone` deste diretório, para que os cinco
% diagramas não divirjam em fonte, cor ou espessura de traço.
%
% fontenc T1 + inputenc utf8 são obrigatórios: sem eles "ç", "ã" e "é"
% saem trocados ou desaparecem silenciosamente do PDF.
\usepackage[T1]{fontenc}
\usepackage[utf8]{inputenc}
\usepackage[brazil]{babel}
% A ordem importa: `amssymb` antes de `newtxmath`, exatamente como em
% `doeapostila.cls`. Invertida, as duas definem `\Bbbk` e a compilação aborta.
\usepackage{amsmath,amssymb}
\usepackage{newtxtext,newtxmath}   % mesma família serifada da classe doeapostila
\usepackage{tikz}
\usetikzlibrary{arrows.meta,calc,decorations.pathreplacing,positioning,
  shapes.geometric,backgrounds}

% Paleta: preto, cinzas e uma única cor de apoio.
\definecolor{azulDoE}{HTML}{1B4E79}
\definecolor{cinzaDoE}{gray}{0.55}
\definecolor{cinzaClaro}{gray}{0.82}
\definecolor{cinzaFundo}{gray}{0.93}

\tikzset{
  % vértice de plano fatorial, com a resposta escrita dentro
  vertice/.style={circle, draw=black, fill=white, line width=0.5pt,
                  inner sep=0pt, minimum size=8.5mm, font=\scriptsize},
  vertdestaque/.style={circle, draw=azulDoE, fill=azulDoE!10, text=azulDoE,
                  line width=0.9pt, inner sep=0pt, minimum size=8.5mm,
                  font=\scriptsize},
  % ponto de planejamento, sem texto dentro
  ponto/.style={circle, draw=black, fill=black, inner sep=0pt,
                minimum size=2.4mm},
  pontovazio/.style={circle, draw=black, fill=white, line width=0.5pt,
                inner sep=0pt, minimum size=2.6mm},
  pontoazul/.style={circle, draw=azulDoE, fill=azulDoE, inner sep=0pt,
                minimum size=2.8mm},
  axial/.style={rectangle, draw=azulDoE, fill=azulDoE, inner sep=0pt,
                minimum size=2.6mm},
  apagado/.style={circle, draw=cinzaClaro, fill=white, line width=0.5pt,
                inner sep=0pt, minimum size=2.6mm},
  aresta/.style={draw=black, line width=0.5pt},
  arestafundo/.style={draw=cinzaClaro, line width=0.5pt},
  eixo/.style={draw=black, line width=0.5pt, -{Latex[length=1.8mm,width=1.4mm]}},
  seta/.style={draw=black, line width=0.6pt, -{Latex[length=2.2mm,width=1.7mm]}},
  setaazul/.style={draw=azulDoE, line width=0.8pt,
                   -{Latex[length=2.2mm,width=1.7mm]}},
  cota/.style={draw=cinzaDoE, line width=0.5pt,
               {Latex[length=1.6mm]}-{Latex[length=1.6mm]}},
  % Corpos escolhidos para casar com as figuras de matplotlib, que usam
  % 8 pt nos ticks e 9 pt nos rótulos de eixo. Com a base de 11 pt do
  % `standalone` (a mesma da classe doeapostila), \scriptsize = 8 pt e
  % \footnotesize = 9 pt. Abaixo disso a figura deixa de ser legível na
  % página impressa, e é para isso que a folha de prova serve.
  rotulo/.style={font=\footnotesize},
  rotulinho/.style={font=\scriptsize},
  nota/.style={font=\footnotesize, align=left},
  notaazul/.style={font=\footnotesize, align=left, text=azulDoE},
}

\begin{document}
\begin{tikzpicture}[x=1.15cm, y=1.15cm,
  declare function={
    reta(\x)        = 0.32*\x;
    semilargura(\x) = 0.95*sqrt(0.16 + (\x)*(\x)/3.2);
    sup(\x) = reta(\x) + semilargura(\x);
    inf(\x) = reta(\x) - semilargura(\x);
  }]

  \def\xa{-3.55}
  \def\xb{3.55}
  \def\ya{-3.62}
  \def\yb{3.62}
  \def\xexp{2.3}   % limite da faixa efetivamente ensaiada

  % --- zonas de extrapolação ---------------------------------------------
  \fill[cinzaFundo] (\xa,\ya) rectangle (-\xexp,\yb);
  \fill[cinzaFundo] (\xexp,\ya) rectangle (\xb,\yb);
  \node[rotulinho, text=cinzaDoE, align=center] at (-2.92,3.15)
    {extra-\\polação};
  \node[rotulinho, text=cinzaDoE, align=center] at (2.92,3.15)
    {extra-\\polação};

  % --- eixos --------------------------------------------------------------
  \draw[eixo] (\xa,\ya) -- (\xb+0.35,\ya) node[rotulo, right] {$x$};
  \draw[eixo] (\xa,\ya) -- (\xa,\yb+0.35) node[rotulo, above] {$\hat{y}$};

  % --- banda --------------------------------------------------------------
  \draw[azulDoE, line width=0.9pt, smooth, samples=140, domain=\xa:\xb]
        plot (\x, {sup(\x)});
  \draw[azulDoE, line width=0.9pt, smooth, samples=140, domain=\xa:\xb]
        plot (\x, {inf(\x)});

  % --- reta ajustada ------------------------------------------------------
  \draw[black, line width=1pt, domain=\xa:\xb] plot (\x, {reta(\x)});
  \node[rotulo, anchor=north west, rotate=23] at (1.35,{reta(1.35)-0.08})
    {$\hat{y} = \hat{\beta}_0 + \hat{\beta}_1 x$};

  % --- média dos x --------------------------------------------------------
  \draw[cinzaDoE, dashed, line width=0.5pt] (0,\ya) -- (0,{sup(0)});
  \node[rotulinho, below] at (0,\ya-0.06) {$\bar{x}$};

  % --- cota do estrangulamento -------------------------------------------
  \draw[cota] (0,{inf(0)}) -- (0,{sup(0)});
  \node[notaazul, anchor=south east, align=right] at (-0.55,1.55)
    {a banda é mais estreita\\ em $\bar{x}$: é ali que a\\ previsão é mais firme};
  \draw[setaazul] (-0.62,1.42) .. controls (-0.35,0.95) .. (-0.06,{sup(0)+0.05});

  % --- cota da abertura na borda -----------------------------------------
  \draw[cota] (\xb-0.10,{inf(\xb-0.10)}) -- (\xb-0.10,{sup(\xb-0.10)});
  \node[notaazul, anchor=north west, align=left] at (0.55,-2.35)
    {e se abre nas bordas: prever longe\\ do centro custa incerteza};
  \draw[setaazul] (2.72,-2.30) -- (\xb-0.20,{inf(\xb-0.20)-0.12});

  % --- faixa experimental -------------------------------------------------
  \draw[decorate, decoration={brace, amplitude=4pt, mirror}, cinzaDoE,
        line width=0.5pt] (-\xexp,\ya-0.42) -- (\xexp,\ya-0.42);
  \node[rotulo, below] at (0,\ya-0.78) {faixa efetivamente ensaiada};

\end{tikzpicture}
\end{document}
